\documentclass[a4paper,12pt]{article}
\usepackage[italian]{babel}
\usepackage[utf8]{inputenc}
\usepackage[T1]{fontenc}
\usepackage{geometry}
\usepackage{setspace}
\geometry{margin=2.5cm}
\onehalfspacing

\title{Verifica scritta di informatica - Classe 4 AIQ - Compito A}
\author{Prof. Fedeli Massimo}
\date{}
\begin{document}
	\maketitle
	
	\textbf{Durata:} 2 ore \\
	\textbf{Obiettivo:} scrivere interrogazioni SQL (anche nidificate) e formulare interrogazioni in algebra relazionale.
	
	\section*{Schema logico}
	
	\begin{flushleft}
		CLIENTI(id\_cliente PK, cognome, nome, email, telefono, data\_reg ) \\
		MOTO(id\_moto PK, targa , marca, modello, cilindrata , anno , categoria, tariffa\_giornaliera, attiva ) \\
		NOLEGGI(id\_noleggio PK, id\_cliente FK$\rightarrow$CLIENTI, id\_moto FK$\rightarrow$MOTO, data\_inizio, data\_fine , km\_inizio , km\_fine , cauzione ) \\
		PAGAMENTI(id\_pagamento PK, id\_noleggio FK$\rightarrow$NOLEGGI , data\_pagamento , importo , metodo, esito) \\
		MECCANICI(id\_meccanico PK, cognome, nome, livello) \\
		INTERVENTI(id\_intervento PK, id\_moto FK$\rightarrow$MOTO, id\_meccanico FK$\rightarrow$MECCANICI, data\_intervento , tipo, ore\_manodopera , costo\_manodopera\_oraria , note) \\
		RICAMBI(id\_ricambio PK, descrizione, costo\_unitario , giacenza ) \\
		INTERVENTI\_RICAMBI(id\_intervento FK$\rightarrow$INTERVENTI, id\_ricambio FK$\rightarrow$RICAMBI, quantita , PK(id\_intervento,id\_ricambio))
	\end{flushleft}
	
	\section*{Parte A -- Query SQL - 7.5 pt}
	
	Scrivere le seguenti interrogazioni in linguaggio SQL. (Per ottenere il punteggio massimo è sufficiente svolgere 9 esercizi su 10 a scelta dello studente)
	
	\begin{enumerate}
		\item Mostrare tutti i noleggi effettuati dal cliente con id\_cliente = 7 indicando id\_noleggio, targa, data\_inizio e data\_fine, nome e cognome del cliente.
		\item Elencare i noleggi iniziati nel mese di febbraio 2026 mostrando cliente, moto e periodo di noleggio.
		\item Elencare le moto che non sono mai state noleggiate.
		\item Elencare i clienti che hanno noleggiato almeno una moto della categoria ``Touring''.
		\item Elencare i noleggi che non hanno un pagamento con esito positivo.
		\item Per ogni moto calcolare il numero totale di noleggi effettuati.
		\item Per ogni cliente calcolare la spesa totale sostenuta considerando solo i pagamenti con esito positivo.
		\item Elencare le moto il cui costo totale di manutenzione è maggiore del costo medio di manutenzione tra tutte le moto.
		\item Trovare i clienti che hanno noleggiato almeno una volta ogni categoria di moto attiva.
		\item Per ciascun cliente mostrare il noleggio più recente effettuato.
	\end{enumerate}
	
	\section*{Parte B -- Algebra Relazionale 2.5 pt}
	
	Scrivere le seguenti interrogazioni in algebra relazionale. 
	
	\begin{enumerate}
		\item Targhe delle moto attive di categoria ``Touring''.
		\item Cognome e nome dei clienti registrati nel 2026.
		\item Targhe delle moto che compaiono in almeno un noleggio.
		\item Coppie (cliente, targa) relative ai noleggi di febbraio 2026.
	\end{enumerate}
	
\end{document}
